\chapter{Mar.~31 --- Riemann-Roch Theorem}

\section{More on Vanishing of Cohomology}

\begin{theorem}
  If $X$ is a projective variety
  and $\mathcal{F}$ is a
  coherent sheaf on $X$, then
  \[
    H^i(X, \mathcal{F}) = 0
  \]
  for all $i > \dim \Supp(\mathcal{F})$.
\end{theorem}

\begin{proof}
  Set $r = \dim \Supp(\mathcal{F})$ and choose
  hyperplanes $H_0, \dots, H_r$ such that
  \[
    \Supp(\mathcal{F})
    \cap H_0 \cap \dots \cap H_r = \varnothing.
  \]
  Now take $U_i = X \setminus H_i$, which
  are open affines with
  \[
    (U_0 \cup \dots \cup U_r) \cap \Supp(\mathcal{F}) = \Supp(\mathcal{F}).
  \]
  Choose an open affine cover
  $U_{r + 1}, \dots, U_s$ such that
  $U_{r + 1} \cup \cdots \cup U_s = X \setminus \Supp(\mathcal{F})$.
  So
  \[
    H^i(X, \mathcal{F})
    = H^i_{U_0, \dots, U_s}(X, \mathcal{F})
    = H^i_{U_0, \dots, U_r}(X, \mathcal{F})
  \]
  since $\mathcal{F}(U_i) = 0$ for $i > r$.
  Thus $H^i(X, \mathcal{F}) = 0$ for $i > r$.
\end{proof}

\section{Riemann-Roch Theorem}

\begin{remark}
  The \emph{Riemann-Roch problem} is the following:
  Let $C$ be a smooth projective curve and
  $\mathcal{L}$ an invertible sheaf on $C$.
  What is $h^0(C, \mathcal{L}) = \dim \Gamma(C, \mathcal{L})$?
  The following are special cases:
  \begin{itemize}
    \item If $\mathcal{L} = \omega_C$, then
      $h^0(C, \omega_C)$ is the dimension
      of the space of algebraic $1$-forms.
    \item If $\mathcal{L} = \omega_C(D) = \omega_C \otimes \OO_C(D)$,
      then $h^0(C, \omega_C(D))$
      gives the dimension of the space of algebraic $1$-forms
      with certain conditions on poles.
    \item If $\mathcal{L} = \OO_C$, then
      $h^0(C, \OO_C) = 1$.
  \end{itemize}
  We can approach this problem as follows:
  \begin{itemize}
    \item instead compute
      $\chi(C, \mathcal{L}) = h^0(C, \mathcal{L}) - h^1(C, \mathcal{L})$;
    \item using $\Cl(C) \overset{\cong}{\longrightarrow} \Pic(C)$
      given by $[D] \mapsto \OO_C(D)$,
      we may assume that $\mathcal{L} = \OO_C(D)$.
  \end{itemize}
\end{remark}

\begin{theorem}[Riemann-Roch]\label{thm:riemann-roch}
  If $C$ is a smooth projective curve and
  $D$ is a divisor on $C$, then
  \[
    \chi(C, \OO_C(D))
    = \deg D + 1 - g,
  \]
  where $g = h^1(C, \OO_C)$ is the
  \emph{genus} of $C$, and
  $\deg D = \sum_{i = 1}^r a_i$ where
  $D = \sum_{i = 1}^r a_i P_i$.
\end{theorem}

\begin{remark}
  By Serre duality, we can rewrite
  \begin{align*}
    \chi(C, \OO_C(D))
    &= h^0(C, \OO_C(D)) - h^1(C, \OO_C(D)) \\
    &= h^0(C, \OO_C(D)) - h^0(C, \OO_C(-D) \otimes \omega_C) \\
    &= h^0(C, \OO_C(D)) - h^0(C, \OO_C(K_C - D)),
  \end{align*}
  where $K_C$ is the canonical divisor
  of $C$. Riemann showed (1857) that
  \[
    h^0(C, \OO_C(D)) \ge \deg D + 1 - g,
  \]
  and Roch improved this (1865) to
  the full statement of Theorem \ref{thm:riemann-roch}:
  \[
    h^0(C, \OO_C(D))
    - h^0(C, \OO_C(K_C - D))
    = \deg D + 1 - g.
  \]
\end{remark}

\begin{proof}[Proof of Theorem \ref{thm:riemann-roch}]
  We proceed by induction. The base case is
  $D = 0$, and we have
  \[
    \chi(C, \OO_C)
    = h^0(C, \OO_C) - h^1(C, \OO_C)
    = 1 - g = \deg(0) + 1 - g.
  \]
  Now we show the inductive step. It suffices
  to show that for any $P \in C$, the theorem
  holds for $D$ if and only if it holds for
  $D - P$ (then we can repeatedly add or
  subtract points until we get to the desired
  divisor).
  Consider the short exact sequence
  \[
    \begin{tikzcd}
      0 \ar[r] & \OO_C(-P) \ar[r] & \OO_C \ar[r] & \OO_P \ar[r] & 0,
    \end{tikzcd}
  \]
  and apply $- \otimes \OO_C(D)$ to get
  a short exact sequence (note that
  the sheaves above are locally free)
  \[
    \begin{tikzcd}[row sep = small]
      & & & \OO_P \ar[d, equal] \\
      0 \ar[r] & \OO_C(D - P) \ar[r] & \OO_C(D) \ar[r] & \OO_P \otimes \OO_C(D) \ar[r] & 0.
    \end{tikzcd}
  \]
  Then we have that
  \begin{align*}
    \chi(C, \OO_C(D))
    - \chi(C, \OO_C(D - P))
    &= \chi(C, \OO_P)
    = h^0(C, \OO_P) - h^1(C, \OO_P) \\
    &= h^0(C, \OO_P)
    = 1 = \deg D - \deg (D - P),
  \end{align*}
  so Riemann-Roch holds for $D$ if and only if
  it holds for $D - P$.
\end{proof}
